
\documentclass{article}
\usepackage{colortbl}
\usepackage{makecell}
\usepackage{multirow}
\usepackage{supertabular}

\begin{document}

\newcounter{utterance}

\centering \large Interaction Transcript for game `wordle', experiment `medium\_frequency\_words\_no\_clue\_no\_critic', episode 20 with Qwen3.5{-}4B.
\vspace{24pt}

{ \footnotesize  \setcounter{utterance}{1}
\setlength{\tabcolsep}{0pt}
\begin{supertabular}{c@{$\;$}|p{.15\linewidth}@{}p{.15\linewidth}p{.15\linewidth}p{.15\linewidth}p{.15\linewidth}p{.15\linewidth}}
    \# & \multicolumn{2}{c}{Player} && \multicolumn{2}{c}{Game Master} \\
    \hline

    \theutterance \stepcounter{utterance}  
    & & & \multicolumn{4}{p{0.6\linewidth}}{
        \cellcolor[rgb]{0.9,0.9,0.9}{
            \makecell[{{p{\linewidth}}}]{
                \texttt{\tiny{[P1$\langle$GM]}}
                \texttt{You are a language wizard who likes to guess words by using the given rules.} \\
\\ 
\texttt{Welcome to Wordle! You have six attempts to guess the target word, a valid English word of five lowercase letters (a{-}z). Please use the tags "explanation:" and "guess:" to provide a concise explanation for each guess.} \\
\\ 
\texttt{For instance, if your guess is "apple", your response should be} \\
\texttt{explanation: this is a common five{-}letter English word, and I am starting my guess with this word.} \\
\texttt{guess: apple} \\
\\ 
\texttt{After each guess, your answer will be validated, and you will receive feedback indicating which letters are correct (green), which letters are correct but in the wrong position (yellow), and which letters are incorrect (red). This feedback can be useful in determining which letters to include or exclude in your next guess.} \\
\\ 
\texttt{For example, the feedback for "apple" might be:} \\
\texttt{guess\_feedback: a<yellow> p<yellow> p<green> l<yellow> e<red>} \\
\\ 
\texttt{The explanation should contain details about how the guess\_feedback is used to arrive at a new guess.} \\
\\ 
\texttt{Do not generate any other text.} \\
\\ 
\texttt{Let's begin with your first guess.} \\
            }
        }
    }
     \\ \\

    \theutterance \stepcounter{utterance}  
    & \multicolumn{4}{p{0.6\linewidth}}{
        \cellcolor[rgb]{0.561,0.737,0.561}{
            \makecell[{{p{\linewidth}}}]{
                \texttt{\tiny{[P1$\rangle$GM]}}
                \texttt{explanation: I will start with a very common five{-}letter word that contains a variety of frequent letters to maximize information gain.} \\
\texttt{guess: apple} \\
            }
        }
    }
    & & \\ \\

    \theutterance \stepcounter{utterance}  
    & & & \multicolumn{2}{p{0.3\linewidth}}{
        \cellcolor[rgb]{0.9,0.9,0.9}{
            \makecell[{{p{\linewidth}}}]{
                \texttt{\tiny{[GM$|$GM]}}
                \texttt{attempts = 1} \\
\texttt{target\_word = glove} \\
\texttt{guess = apple} \\
\texttt{guess\_feedback = a<red> p<red> p<red> l<yellow> e<green>} \\
            }
        }
    }
    & & \\ \\

    \theutterance \stepcounter{utterance}  
    & & & \multicolumn{4}{p{0.6\linewidth}}{
        \cellcolor[rgb]{0.9,0.9,0.9}{
            \makecell[{{p{\linewidth}}}]{
                \texttt{\tiny{[P1$\langle$GM]}}
                \texttt{guess\_feedback: a<red> p<red> p<red> l<yellow> e<green>} \\
\\ 
\texttt{Provide your response only in this format.} \\
\texttt{explanation: details} \\
\texttt{guess: word} \\
            }
        }
    }
     \\ \\

    \theutterance \stepcounter{utterance}  
    & \multicolumn{4}{p{0.6\linewidth}}{
        \cellcolor[rgb]{0.561,0.737,0.561}{
            \makecell[{{p{\linewidth}}}]{
                \texttt{\tiny{[P1$\rangle$GM]}}
                \texttt{explanation: The feedback "e<green>" and "l<yellow>" tells me the target word ends in "e" and contains an "l" somewhere before the last letter, but not in the fourth position. The letters "a", "p", and the second "p" are incorrect. I need to form a new word ending in "e" that includes an "l" in a different position. "olive" fits these criteria: it ends in "e", contains "l", and uses common letters.} \\
\texttt{guess: olive} \\
            }
        }
    }
    & & \\ \\

    \theutterance \stepcounter{utterance}  
    & & & \multicolumn{2}{p{0.3\linewidth}}{
        \cellcolor[rgb]{0.9,0.9,0.9}{
            \makecell[{{p{\linewidth}}}]{
                \texttt{\tiny{[GM$|$GM]}}
                \texttt{attempts = 2} \\
\texttt{target\_word = glove} \\
\texttt{guess = olive} \\
\texttt{guess\_feedback = o<yellow> l<green> i<red> v<green> e<green>} \\
            }
        }
    }
    & & \\ \\

    \theutterance \stepcounter{utterance}  
    & & & \multicolumn{4}{p{0.6\linewidth}}{
        \cellcolor[rgb]{0.9,0.9,0.9}{
            \makecell[{{p{\linewidth}}}]{
                \texttt{\tiny{[P1$\langle$GM]}}
                \texttt{guess\_feedback: o<yellow> l<green> i<red> v<green> e<green>} \\
\\ 
\texttt{Provide your response only in this format.} \\
\texttt{explanation: details} \\
\texttt{guess: word} \\
            }
        }
    }
     \\ \\

    \theutterance \stepcounter{utterance}  
    & \multicolumn{4}{p{0.6\linewidth}}{
        \cellcolor[rgb]{0.561,0.737,0.561}{
            \makecell[{{p{\linewidth}}}]{
                \texttt{\tiny{[P1$\rangle$GM]}}
                \texttt{explanation: The feedback "o<yellow> l<green> i<red> v<green> e<green>" confirms the pattern \_l\_v\_e. The letter 'o' is in the word but not in the first position, and 'i' is incorrect. Since 'o' is yellow in the guess "olive" (where it is in the first position), the first letter must be a different letter that forms a valid word with the pattern \_l\_v\_e. The word "olive" itself is the only common word fitting \_o\_l\_v\_e, but the feedback indicates the first letter is not 'o'. However, re{-}evaluating the feedback: 'o' is yellow means it exists in the word but not in the position guessed (1st). 'l' is green (2nd), 'v' is green (4th), 'e' is green (5th). The pattern is \_l\_v\_e. The letter 'o' is present but not first. This implies the first letter is 'o' but the feedback says it's yellow? No, if 'o' is yellow in "olive", it means 'o' is in the word but not at index 0. But 'o' is at index 0 in "olive". This is a contradiction unless the word is not "olive". Wait, the feedback says 'o' is yellow. That means 'o' is in the word, but not at position 0. But in the guess "olive", 'o' is at position 0. Therefore, the word cannot be "olive". The word must have 'o' somewhere else. But the pattern is \_l\_v\_e. The only slot for 'o' is position 0. If 'o' is at position 0, it should be green. Since it is yellow, 'o' is not at position 0. But the pattern \_l\_v\_e requires a letter at position 0. If that letter is not 'o', then 'o' must be elsewhere. But there is no other slot. This suggests the word is not "olive". Let's reconsider the feedback. Maybe the word is "olive" and the feedback logic in the prompt example was hypothetical or I am misinterpreting "yellow". If 'o' is yellow, it means 'o' is in the word but not in the first position. But 'o' is in the first position in the guess. This is a logical impossibility for a standard Wordle game unless the word is not "olive". However, assuming the feedback is correct, the word must contain 'o' but not at the start. The pattern is \_l\_v\_e. The only way 'o' is not at the start is if the first letter is something else, say 'x', and 'o' is... where? There is no other slot. This implies the word is actually "olive"} \\
            }
        }
    }
    & & \\ \\

    \theutterance \stepcounter{utterance}  
    & & & \multicolumn{2}{p{0.3\linewidth}}{
        \cellcolor[rgb]{0.9,0.9,0.9}{
            \makecell[{{p{\linewidth}}}]{
                \texttt{\tiny{[GM$|$GM]}}
                \texttt{Error: The response should contain the 'guess:' keyword exactly once.} \\
            }
        }
    }
    & & \\ \\

    \theutterance \stepcounter{utterance}  
    & & & \multicolumn{2}{p{0.3\linewidth}}{
        \cellcolor[rgb]{0.9,0.9,0.9}{
            \makecell[{{p{\linewidth}}}]{
                \texttt{\tiny{[GM$|$GM]}}
                \texttt{game\_result = ABORT} \\
            }
        }
    }
    & & \\ \\

\end{supertabular}
}

\end{document}
